% ICCV 2025 Paper Template

\documentclass[10pt,twocolumn,letterpaper]{article}

%%%%%%%%% PAPER TYPE  - PLEASE UPDATE FOR FINAL VERSION
% \usepackage{iccv}              % To produce the CAMERA-READY version
%\usepackage[review]{iccv}      % To produce the REVIEW version
\usepackage[pagenumbers]{iccv} % To force page numbers, e.g. for an arXiv version
\usepackage{multirow}
\usepackage{xcolor} % 加载颜色包
\usepackage{amssymb} % 加载符号包
% Import additional packages in the preamble file, before hyperref
%
% --- inline annotations
%
\newcommand{\red}[1]{{\color{red}#1}}
\newcommand{\todo}[1]{{\color{red}#1}}
\newcommand{\TODO}[1]{\textbf{\color{red}[TODO: #1]}}
% --- disable by uncommenting  
% \renewcommand{\TODO}[1]{}
% \renewcommand{\todo}[1]{#1}



% It is strongly recommended to use hyperref, especially for the review version.
% hyperref with option pagebackref eases the reviewers' job.
% Please disable hyperref *only* if you encounter grave issues, 
% e.g. with the file validation for the camera-ready version.
%
% If you comment hyperref and then uncomment it, you should delete *.aux before re-running LaTeX.
% (Or just hit 'q' on the first LaTeX run, let it finish, and you should be clear).
\definecolor{iccvblue}{rgb}{0.21,0.49,0.74}
\usepackage[pagebackref,breaklinks,colorlinks,allcolors=iccvblue]{hyperref}

%%%%%%%%% PAPER ID  - PLEASE UPDATE
\def\paperID{9236} % *** Enter the Paper ID here
\def\confName{ICCV}
\def\confYear{2025}

%%%%%%%%% TITLE - PLEASE UPDATE
\title{Supplement Material for MuseTalk: Real-Time High-Fidelity Video Dubbing via Spatio-Temporal Sampling}

\begin{document}
\maketitle


\section{Data Filtering}
% In this section, we provide a detailed description of the datasets used during the training of MuseTalk, including the basic information of our proprietary dataset and the details of the data filtering pipeline.

% \paragraph{Proprietary Dataset.}
% Our proprietary dataset consists of 2,011 video clips with the duration around 5 seconds.  
% The majority of these clips are Chinese broadcast videos, featuring a number of Asian identities. 
% This is intended to compensate for the over-representation of Caucasian identities in our model, as both VFHQ~\cite{xie2022vfhq} and HDTF~\cite{zhang2021flow} datasets predominantly contain Caucasian identities.
We've discovered that the audio-visual consistency of the original video is crucial for model training of video-dubbing task, making data filtering an essential step. 
For instance, the VFHQ dataset contains numerous interview videos where the audio originates from someone off-camera, while the person on-screen remains silent. This type of ``dirty" data can significantly impair the effectiveness of video dubbing methods.

Our filtering approach involves calculating scores based on SyncNet, and subsequently eliminating data with low confidence and high offset. For a comparison of the datasets before and after filtering, please refer to~\cref{tab:dataset_statistics}.



\begin{table}[t]
\begin{tabular}{lccc}
\toprule
\textbf{Dataset}   & \textbf{Original Dataset}& \textbf{Filtered Dataset}\\
\midrule
HDTF~\cite{zhang2021flow}   & 15.75 & 15.32 \\
VFHQ~\cite{xie2022vfhq}   & 18.36 &  8.43 \\
\bottomrule
\end{tabular}
\caption{The variation in the duration of the dataset before and after filtering. The units in the table are measured in hours.}
\label{tab:dataset_statistics}
\end{table}

\section{Ablation Study for Multi Stage Training}
As mentioned in the Section 3 in main text, optimizing \(\mathcal{L}_{\text{adv}}\) and \(\mathcal{L}_{\text{sync}}\) simultaneously in a single training stage for a less capable (randomly initialized) model leads to severe training instability. We illustrate this issue through qualitative results. As shown in~\cref{fig:fig_gan}, single-stage training introduces significant artifacts, including white spots at the corners of the mouth, missing mouth corners, and jagged teeth.

In contrast, the two-stage training strategy employed by MuseTalk first employs a more moderate training regimen in first stage to equip the model with initial face inpainting capabilities. 
This approach stabilizes the model, enabling it to achieve higher-quality results during the adversarial training in the scoond stage. 
These findings demonstrate the necessity of our proposed two-stage training strategy.
\begin{figure}[t]
\centering
\includegraphics[width=0.8\linewidth]{imageSup/supp_multi_stage.pdf}
\caption{
Artifacts observed in the generation results from single-stage training.}
\label{fig:fig_gan}
\end{figure}

\paragraph{ID Diversity in Different Stage.}
The diagram in Fig.~\ref{fig:fig_stage12} demonstrates our application of different data sampling strategies at various training stages.

In the first stage, our goal is to expose the model to a wide array of unique identities. 
This exposure enables the model to learn a diverse range of facial details, thereby enhancing its capacity for facial abstraction.
To accomplish this, we employ a larger batch size and sample only one frame per video (identity).

In the second stage, we shift our focus towards optimizing the consistency of lip movements for a specific identity.
This stage is pivotal as it ensures the model's proficiency in accurately synchronizing lip movements with the corresponding audio, a critical aspect of video dubbing.

During the initial stage, we can configure the batch size per GPU to 32.
However, in the second stage, due to the necessity of concurrently sampling N frames to compute the sync loss (where N equals 16 in our case), the batch size per GPU is set to 2.

\begin{figure}[tbp]
\centering
\includegraphics[width=\linewidth]{imageSup/stage12_sampling.pdf}
\caption{
Data sampling strategies for different training stages.
}
\label{fig:fig_stage12}
\end{figure}



\section{Details for Spatial-Temporal Sampling}

\paragraph{Visualization for IFS.}

\begin{figure}[t]
\centering
\includegraphics[width=0.8\linewidth]{imageSup/random_sampling.pdf}
\caption{
Comparison between random sampling and Informative Frame Sampling (IFS).}
\label{fig:fig_random}
\end{figure}
We visualize the sampled data using Informative Frame Sampling (IFS) in Fig.~\ref{fig:fig_random}. 
As shown, IFS is capable of selecting \( I^{t}_{ref} \) frames that are consistent in pose but inconsistent in lip movement with the ground-truth frame \( I^{t}_{gt} \). 
In contrast, random sampling may yield input data with large pose variations but similar lip movements compared to \( I^{t}_{gt} \).


\begin{figure*}[tbp]
\centering
\includegraphics[width=\linewidth]{imageSup/lip_crop.pdf}
\caption{
Visualization of the differences in processing methods for different lip regions.
}
\label{fig:fig_lip}
\end{figure*}


\section{Details for Loss Function}

\paragraph{SyncNet Loss.}
The commonly used SyncNet model for evaluating audio-driven video generation is derived from Wav2Lip~\cite{prajwal2020lip} training, which is trained on very low-resolution images ($96\times96$ pixels). 
After reviewing prior works~\cite{tian2024emo, peng2024synctalk}, we found that this model tends to produce higher lip-sync-error confidence (LSE-C)~\citep{prajwal2020lip} for low-resolution generation methods, which does not strictly reflect the lip-sync quality of the model.

Consequently, we sought a more effective audio-visual synchronization guidance model.
Recently, LatentSync~\cite{li2024latentsync} proposed a more effective SyncNet training strategy, training a SyncNet with $N = 16$ to handle $256\times256$ image inputs.
This SyncNet has approximately ten times the number of parameters compared to the model provided by Wav2Lip. Upon manual inspection, we found that this model provides more accurate audio-visual synchornize scores, which we integrated into our training pipeline.

It is crucial to note that MuseTalk's contribution lies not in providing a superior SyncNet or a more innovative loss function, but rather in harmonizing these classical losses to achieve higher-quality generation.



\paragraph{Adversarial Loss.}
During the training process, we utilize two distinct discriminators: \(\mathcal D_{face}\), which concentrates on the entire face, and \(\mathcal D_{lip}\), which specifically targets the lip region.  
The adversarial loss is incorporated using the WGAN (Wasserstein Generative Adversarial Network) framework~\cite{arjovsky2017wasserstein}, offering a more stable training experience compared to conventional GANs.
For the generator, the adversarial loss's objective is to deceive both discriminators into accepting the generated images as authentic. Specifically, the generator's goal is to minimize:
\begin{equation}
\mathcal{L}_{adv} = \mathcal{L}_{adv,face}+\mathcal{L}_{adv,lip},
\end{equation}
where 
\begin{equation}
\mathcal{L}_{adv,face} = -\mathbb{E}_{A_{mel}^{t},I^{t}_{ref}} [D_{face}(I^{t}_{o})],
\end{equation}
\begin{equation}
\mathcal{L}_{adv,lip} = -\mathbb{E}_{A_{mel}^{t},I^{t}_{ref}} [D_{lip}(I^{t}_{lip})].
\end{equation}
For the discriminators, their objectives are to differentiate between real and generated samples. 
The loss functions for the discriminators are defined as:
\begin{equation}
\mathcal{L}_{dis} = \mathcal{L}_{dis,face}+\mathcal{L}_{dis,lip},
\end{equation}

where
\begin{align}
\mathcal{L}_{dis,face} & = \mathbb{E}_{I^{t}_{real,face}}[D_{face}(I^{t}_{real,face})] 
\nonumber
\\
& + \mathbb{E}_{A_{mel}^{t},I^{t}_{ref}} [D_{face}(I^{t}_{o})]
\end{align}

\begin{align}
\mathcal{L}_{dis,lip} & = \mathbb{E}_{I^{t}_{real,lip}}[D_{lip}(I^{t}_{real,lip})] 
\nonumber
\\
& + \mathbb{E}_{A_{mel}^{t},I^{t}_{ref}} [D_{lip}(I^{t}_{lip})]
\end{align}

The discriminator operates on images of a fixed resolution. 
For \(\mathcal D_{face}\), we've matched the input region to that of MuseTalk.  
However, for \(\mathcal D_{lip}\), the lip region's size varies from frame to frame. 

We explored numerous strategies to handle this fluctuation and found that simply resizing the lip region could lead to distorted mouth shapes. This distortion could negatively impact the synchronization of audio and visuals, and compromise the overall realism.
To tackle this problem, we devised the expansion strategy. This method entails pinpointing the longer side of the mouth based on its landmarks, and then using half of this length to define the shorter side. We subsequently expand outwards to set the upper and lower breakpoints in a symmetrical manner. Following this expansion, the area is resized to a standard region, thereby safeguarding the preservation of the original mouth opening size. 
As illustrated in Fig.~\ref{fig:fig_lip}, the expansion operation guarantees that the lip region maintains its natural proportions. This results in superior mouth shape generation and improved audiovisual consistency.


\section{Blending based on Face Parsing}
\begin{figure*}[tbp]
\centering
\includegraphics[width=\linewidth]{imageSup/blending.pdf}
\caption{
Illustration of the blending process used in MuseTalk.}
\label{fig:fig_blending}
\end{figure*}
We recognize that minor seams might be visible at the intersection of the generated region and the original video. To address this, we utilize a blending strategy based on face parsing.

The procedure starts by enlarging the bounding box of the identified face in the source image to create a square region. Within this area, we implement face parsing methods to generate a multi-class label map \footnote{https://github.com/zllrunning/face-parsing.PyTorch}. This map offers semantic segmentation of facial components such as eyes, mouth, nose, and so on.

Using these labels, we preserve the lower half of the face, excluding the nasal region. This choice aids in maintaining natural facial features while circumventing intricate transitions around the nose. Furthermore, we apply a Gaussian blur to the boundaries of the chosen region to ensure seamless transitions during the blending process, thereby creating a blending mask.

With this mask, we seamlessly integrate the generated region from MuseTalk with the original video content. The blending mask governs the contribution from each source.

For real-time applications, we precalculate the blending mask during the preprocessing phase as it relies solely on the original video content. This optimization considerably minimizes computational overhead during runtime, facilitating efficient processing for real-world deployment.

Figure~\ref{fig:fig_blending} depicts the entire blending process, showcasing how the proposed method delivers visually consistent results while maintaining computational efficiency.


\section{Ethical Considerations}
Technologies like MuseTalk, which facilitate video dubbing, contribute significantly to advancements in multimedia communication and entertainment. However, they also raise potential ethical concerns, particularly regarding their misuse for creating deepfakes and other forms of deceptive content. The synthesized results from MuseTalk, while highly realistic, may still display certain visual artifacts. These artifacts could serve as indicators for detecting deepfakes, thereby providing a layer of transparency and assisting in the identification of manipulated content. Furthermore, the development and deployment of such technologies should be accompanied by robust measures to prevent unauthorized use and ensure that the generated content is used ethically and legally.

In summary, while MuseTalk and similar technologies provide powerful tools for video dubbing and digital content creation, their use must be approached with caution and responsibility. It is essential to implement necessary safeguards to mitigate potential misuse and uphold ethical standards in the digital domain.

\section{Limitation and Future work}
While MuseTalk demonstrates a notable improvement in face region resolution ($256\times256$) compared to other state-of-the-art methods, it has yet to reach its full resolution potential. Additionally, certain facial details—such as mustaches, lip shape, and color—are not always well-preserved, which can affect identity consistency. Lastly, occasional jitter is observed due to the single-frame generation process, compromising smoothness.

To address these limitations, future work will focus on incorporating higher-quality training data and integrating a temporal module to reduce jitter and ensure smoother transitions. These enhancements aim to improve both resolution and overall visual consistency. Moreover, incorporating super-resolution models like GFP-GAN~\cite{wang2021towards} as a post-processing step could further elevate output quality in real-world applications.

% \subsection{Performance of the open-sourced weights on HDTF dataset}
% In this appendix, we present the performance metrics for the open-sourced versions of our models, MuseTalk 1.0 and MuseTalk 2.0.
% MuseTalk 1.0 is developed using a training strategy that exclusively relies on L1 loss.
% In contrast, MuseTalk 2.0 adopts the innovative training technique outlined in this paper.

% To ensure a fair comparison with existing open-sourced methods, the models described in the main text are trained solely on the HDTF dataset. We found that models trained only on the HDTF dataset tend to exhibit color discrepancy issues. Incorporating a small amount of proprietary data (talking face videos from Asian individuals) significantly mitigates this problem but results in a slight decrease in lip-sync accuracy. 
% In practical applications, the improvement in color consistency and facial coherence is more noticeable than the loss in lip-sync accuracy. Both versions of our open-sourced models are trained on the "HDTF+proprietary data" dataset. The HDTF dataset encompasses 16 hours of video footage, while the proprietary dataset includes an additional 3 hours of video content. 

% \begin{table}[t]
% \caption{Performance Metrics for HTDF.}
% \label{performance-metrics-htdf}
% \begin{center}
% \begin{tabular}{lccc}
% \midrule
% % \cmidrule(lr){2-4} \cmidrule(lr){5-7}
% & \multicolumn{1}{c}{FID} 
% & \multicolumn{1}{c}{CSIM$\uparrow$} 
% & \multicolumn{1}{c}{LSE-C$\uparrow$} 
% \\
% \hline 
% Wav2Lip \citep{prajwal2020lip} & 11.21 & 0.8184 & 7.46 \\
% VideoRetalking\citep{cheng2022videoretalking} & 10.93 & 0.7989 & \bf{7.7}  \\
% TalkLip \citep{wang2023seeing} & 17.61 & 0.7898 & 2.29  \\
% DI-Net \citep{zhang2023dinet} & 7.27 & 0.8154 & 6.17    \\
% \midrule
% Ground Truth & 0.00 & 1.0 & 7.81  \\
% MuseTalk 2.0 (HDTF) & 6.43 & 0.8225 & 6.53  \\
% \midrule
% MuseTalk 1.0 (HDTF+proprietary data) & 6.38 & 0.8204 & 2.91   \\
% MuseTalk 2.0 (HDTF+proprietary data) & \bf{5.93} & \bf{0.8415} & 4.04  \\
% \bottomrule
% \end{tabular}%
% \end{center}
% \end{table}
{
    \small
    \bibliographystyle{ieeenat_fullname}
    \bibliography{ICCV2025-Author-Kit-Feb/supplement_material}
}
\end{document}
